\subsection{二直線追加Segre射影と二平方の外側交換子}
\begin{proposition}\label{prop:cubic-CY-one-segre-split}
一交換子の二直線追加Segre射影
$(U_0,U)=(\langle ac,ad,bc,bd\rangle,\langle ac,ad,bc,bd,ab\rangle)$
を持つH3包絡は、既分類のskew包絡の場合に限られる。
また $U=U_0+k(a^2+c^2)$ の外側交換子枝はH3を与えない。
\end{proposition}
\begin{proof}
固定出典\cite{AS4CYpinned}の二直線追加射影の全必要被覆
（内部、全等方、非全等方）と二平方の外側交換子の必要表示を用いる。
非全等方の $Q_{p,q}$ は定義からskew多項式環であり、
既存の許容切断分類を適用する。

元関係からの次数三・四の商射影と中心条件を、
命題\ref{prop:cubic-CY-one-segre-square-empty}と同じ全因子特殊化で検証する。
公開検証器\\
\texttt{verify\_cubic\_CY\_one\_segre\_extensions.py}\quad
は以下の四範囲を扱う。
\begin{itemize}
\item 内部枝の $H_s$ 以外の八表示・族は、全許容値で中心三次が零である。
\item 全等方の $T_0,T_1,T_2$ に中心三次が残るのは
$16k^2+1=0$ 又は $48k^2+1=0$ の場合だけであり、
その中心三次は可換な $c,d$ の二元三次式の空間に含まれる。
\item 非全等方の $C_r,E$ の中心三次は全値で零である。
$D_\rho$ に中心三次が残るのは $\rho^2+\rho+1=0$ の場合だけで、中心三次は $kc^3$ である。
\item 二平方で外側交換子を持つ九表示・族は、全許容値で中心三次が零である。
\end{itemize}
上の二元三次と純三乗は整域切断を与えない。

残る内部枝 $H_s(h)$ では、$a,b$ が可換で
\[
 \begin{gathered}
 ca=qac,\quad cb=q^{-1}bc,\quad da=q^{-1}ad,\quad db=qbd,\\
 dc=cd+\kappa f(a,b),\quad q=(4h+1)/(4h-1),\quad \kappa\ne0 .
 \end{gathered}
\]
$c,d$ の個数の濾過の随伴次数環は、最後の補正を零にしたskew多項式環である。
これは関係からの全射とAS4 Hilbert級数の一致による。
その中心三次単項式 $a^ib^jc^kd^l$ は $q^{i-j}=q^{k-l}=1$ を満たす。
$q=1,-1$ は許容域に現れない。
$i+j+k+l=3$ なので、$q$ が原始三乗根のときの純三乗四つだけが可能である。

原始三乗根の場合、補正付きの環でも $a,b$ との可換性は $k-l\equiv0\pmod3$ を強制する。
$c,d$ の濾過次数が3の項は従って $c^3,d^3$ のみである。
$c^3,d^3$ の交換子の濾過次数は高々2なので、
次数2の補正 $(Aa+Bb)cd$ があれば、$c$ との交換子の濾過次数3から $A=B=0$ となる。
残った次数0の項は $a^3,b^3$ の線型結合に限られる。
$[d,c^3]$ と $[c,d^3]$ に出る二次式は、$f$ の三つの対角変換の和であり、
$ab$ 成分だけが三倍で残る。従って $f$ の $ab$ 成分が非零なら
中心三次は $a^3,b^3$ に限られ、整域切断はない。
内部被覆の表で残るのは $S_F$ の $f=a^2+b^2$ だけである。
この最後の枝には次の補題を用いる。
\end{proof}

\begin{lemma}\label{lem:cubic-CY-two-square-weyl-split}
上の交換関係で $f=a^2+b^2$、$q$ が原始三乗根なら、
中心三次のどの非零切断も整域を与えない。
\end{lemma}
\begin{proof}
$G=Aa^3+Bb^3+Cc^3+Dd^3$ の切断が整域であると仮定する。
$C=D=0$ は二元三次の積分解で排除される。
$a$ を逆元化した次数零部分で $x=ca^{-1},y=da^{-1},z=ba^{-1}$ とおくと、
\[
 xz=qzx,\qquad yz=q^{-1}zy,\qquad yx=q^2xy+q\kappa(1+z^2).
\]
$L(z)=\alpha+\beta z^2$ を
$L(q^{-1}z)-q^2L(z)=q\kappa(1+z^2)$ で定める。
両係数 $\alpha,\beta$ は非零である。$u=xy-L(z)$ とおけば
$ux=q^2xu$, $uy=qyu$, $uz=zu$。
$Z=z^3$ を逆元化して $v=uz^{-1}$ とすると、
$v$ は $x,y,z$ と可換である。
$X=x^3,Y=y^3$ と書くと
\[
 XY=\alpha^3+\beta^3 Z^2+(v^3-3\alpha\beta v)Z,qquad
 A+BZ+CX+DY=0 .
\]
第一式は $xy=\alpha+vz+\beta z^2$ の三共役の積である。

$K=k(X,Z,v)$ を、局所化した部分環に含まれる中央の体とする。
$Y$ も $K$ に属し、$x,z$ はこの体上で三次巡回代数を生成する。
実際 $x^3=X,z^3=Z$ と非自明な交換式があるので、
対応する三次シンボル代数からの単位的写像は単射である。
従ってこのシンボルが分裂すれば整域性に矛盾する。
もし $C,D$ のうち一つが零なら、例えば $C=0$ のとき
$Y=-(A+BZ)/D$ は $K(z)/K$ のノルムとなる。$z$ が中央でないため $Z$ は $K$ の立方でなく、
$y,z$ で同じ代数を生成する。$x=(u+L(z))y^{-1}$ を使えばよい。
従ってこの場合は分裂する。$D=0$ も同様である。

以下 $CD\ne0$ とする。$K$ の $k$ 上の超越次数が高々1ならTsenの定理
\cite[Tag 03RF]{StacksTsenBrauer}で直ちに分裂する。
それ以外では $v$ は超越的で、$K$ は $F=k(v)$ 上の射影二次曲線
\[
 C X_0^2+B X_0 Z_0+A X_0 W+
 D\beta^3 Z_0^2+D(v^3-3\alpha\beta v)Z_0W+D\alpha^3 W^2=0
\]
の函数体である。ここで $X=X_0/W,Z=Z_0/W$。
行列式の $v^6$ 係数は $-CD^2/4\ne0$ なので滑らかである。
無限遠 $W=0$ には $k$ 上の点があり、従って曲線は $F$ 上の射影直線である。

この曲線上で三次シンボル $(X,Z)$ の剰余を調べる。
両函数が単元なら剰余は自明である。零・極のある点は以下で尽くされる。
\begin{itemize}
\item $X=0$ の点では $Z$ は非零かつ有限であり、剰余体で
$\alpha^3+\beta^3 Z^2+(v^3-3\alpha\beta v)Z=0$ となる。
$Z$ が立方でなければ、これは三次体拡大における非零元
$\alpha+v z+\beta z^2$ のノルムが零であることを意味するので矛盾する。
従って $Z$ の剰余は立方である。
\item $Z=0$ の点では、$CX^2+AX+D\alpha^3=0$ により、
$X$ の剰余は $k^*$ に属し立方である。
\item $W=0$ の点では $X_0,Z_0$ は共に非零である。
実際 $C,D\beta^3$ は非零であり、無限遠方程式の両端係数である。
従って $X,Z$ の極位数は等しく、比の剰余は $k^*$ に属するので、
シンボルの剰余は自明である。接する場合の重複も同じ議論で扱う。
\end{itemize}
Faddeevの完全列\cite[Theorem 6.4.5]{GilleSzamuelyFaddeev}により、
このシンボル類は $\operatorname{Br}(F)$ から来る。
$F=k(v)$ のBrauer群はTsenの定理で零である。
従って最後の枝でも分裂し、整域性に矛盾する。
\end{proof}

\begin{proposition}\label{prop:cubic-CY-one-segre-two-squares-empty}
対称射影が $U_0+k(a^2+c^2)$、$U_0=\langle ac,ad,bc,bd\rangle$ である
一交換子CY包絡はH3を与えない。
\end{proposition}
\begin{proof}
外側交換子は命題\ref{prop:cubic-CY-one-segre-split}で排除した。
固定出典\cite{AS4CYpinned}の二平方追加型の全必要被覆により、
残る交差交換子は有限チャート $(p,t)$ と無限遠チャート $e$ に尽くされる。
許容域はそれぞれ
\[
 t F_-F_+\ne0,\quad F_\pm=p(t\pm1)^2+2,
 \qquad (128e-1)(128e+1)\ne0 .
\]
元トップの係数ブロックを復元し、関係空間を張る六列を定数小行列式で確認する。
公開検証器\\
\texttt{verify\_cubic\_CY\_one\_segre\_cross.py}\quad
は、この六列から次数別商と中心行列を再構成する。

無限遠チャートでは、$e=0$ のみを別計算し、いずれも中心行列の階数は20である。
有限チャートでは、疎な商射影と中心の一つの20次小行列式の分子・分母を因数分解する。
許容域内でこの検査から外れるのは次の五曲線だけである。
\[
 p=0,\qquad p=-2,\qquad
 (t+1)p+2=0,\qquad (t+1/2)p+1=0,\qquad (t-1)p-2=0.
\]
各曲線で $p$ を $t$ の有理式に戻す分母が零なら元の曲線式に矛盾する。
従ってこの置換で全点を覆う。
元六関係から商と中心行列を曲線ごとに再計算すると、
小行列式・商射影の例外因子は、全て元の許容域の外にあることが分かる。
従って各曲線の全許容値で階数20となり、全中心三次は零である。
\end{proof}
